\section{Related Work}
\label{sec:related}

\subsection{Machine Learning for Credit Scoring}
Credit scoring has long been treated as a binary classification problem: predict whether an applicant will have
difficulty repaying. A large benchmark of classifiers on several credit scoring datasets found that ensemble methods,
including random forests, are among the strongest performers and that progress should be judged on several data
sets and metrics rather than on accuracy alone~\cite{lessmann2015benchmark}. Gradient-boosted decision trees such as
LightGBM~\cite{ke2017lightgbm} are now widely used for this task, since they train quickly on large tables and need
little preprocessing. Credit data are typically imbalanced, which has motivated resampling methods such as
SMOTE~\cite{chawla2002smote}; with an imbalanced target, threshold-free measures such as the area under the ROC curve
are more informative than accuracy.

The Home Credit Default Risk dataset~\cite{homecredit2018} is a public benchmark with more than 300{,}000
applicants, many of whom have little or no credit history. Shreya and Pathak~\cite{shreya2025base} compared random
forest, XGBoost and LightGBM on this dataset, found LightGBM the strongest, and explained it with SHAP and LIME. Our
work builds on that study. It keeps the model family and the explanation methods, and adds what is needed to turn a
prediction into a decision process: decision bands with officer review, recourse, a fairness analysis and a working
application.

\subsection{Post-Hoc Explanations}
SHAP~\cite{lundberg2017shap} attributes a prediction to the input features using Shapley values, and its tree
variant computes them exactly and efficiently for tree ensembles~\cite{lundberg2020tree}. The attributions add up to
the model's output, which makes them well suited to explaining an individual credit decision. LIME~\cite{ribeiro2016lime}
explains a prediction by sampling perturbed inputs around it and fitting an interpretable surrogate model. Its
explanations depend on how the perturbations are drawn, and both LIME and SHAP explanations can change under small
changes to the input~\cite{alvarezmelis2018robustness}. We therefore measure how well the LIME explanations agree
with SHAP for each applicant, instead of assuming that two explanation methods tell the same story.

How explanations should be evaluated is itself an open question. Doshi-Velez and Kim~\cite{doshivelez2017rigorous}
distinguish evaluation with real users on the real task, evaluation with people on simplified tasks, and functionally
grounded evaluation through measurable proxies. We follow the last approach, with measures of agreement, stability
and readability~\cite{flesch1948readability}, and prepare an expert review based on usability
heuristics~\cite{nielsen1994heuristics}.

\subsection{Counterfactual Explanations and Algorithmic Recourse}
A counterfactual explanation describes the smallest change to an input that would change the model's
decision~\cite{wachter2018counterfactual}. Algorithmic recourse adds the requirement that the change be actionable:
it should only alter features that a person can change, in directions and amounts that are
feasible~\cite{ustun2019recourse}. Ustun~\emph{et al.}~\cite{ustun2019recourse} computed such recourse for linear
classifiers with integer programming. DiCE~\cite{mothilal2020dice} generates several diverse counterfactuals for any
differentiable or black-box model, trading off proximity to the input against diversity among the options. Karimi
\emph{et al.}~\cite{karimi2022survey} survey the area and stress that recourse must account for the cost to the
person and for the plausibility of the suggested profile.

Our recourse engine shares these goals but is designed around credit decisions. Its levers are the ones an applicant
can act on (the loan amount, documented income, the repayment period and the passage of time), immutable attributes
are never changed, and options are offered as a path through the decision bands rather than as a single flip. We
compare it with DiCE and with random search on the same applicants.

\subsection{Explainable Credit Decision Systems}
Several recent studies combine these ideas for credit. Nallakaruppan~\emph{et al.}~\cite{nallakaruppan2024risks}
apply SHAP and LIME to decision trees and random forests for peer-to-peer lending and frame the result as decision
support, but their explanations are attributions only. Dastile and Celik~\cite{dastile2024counterfactual} generate
credit-scoring counterfactuals with genetic algorithms, particle swarm optimization and Bayesian optimization, and
evaluate their plausibility, validity and sparsity; Nguyen Thanh~\emph{et al.}~\cite{nguyen2025milp} obtain plausible
counterfactuals for credit approval with mixed-integer programming and local outlier factor modeling. In these works
the counterfactual algorithm is the whole contribution. Kapoor~\emph{et al.}~\cite{kapoor2026fico} combine TabNet,
LightGBM, SHAP and counterfactual explanations with a fairness perspective on the FICO dataset, and Popoola and
Sheppard~\cite{popoola2026fair} study whether counterfactual explanations are consistent across groups as a measure of
procedural fairness.

CreditXplain differs from these studies in three ways. It integrates the model, both explanation methods, recourse and
a fairness analysis into one decision process with officer review and an applicant portal, rather than evaluating a
single component. Its recourse is evaluated for what happens when the model is retrained, which none of these studies
measure. And it treats part of the fairness question as an eligibility policy outside the model, supported by an
experiment showing that removing age from the model does not narrow the age gap.

\subsection{Robustness of Recourse}
A recourse option is only useful if it still works when the applicant acts on it. Rawal~\emph{et al.}~\cite{rawal2020shift}
showed that recourse computed for one model is often invalid after the model is updated on shifted data, and
Pawelczyk~\emph{et al.}~\cite{pawelczyk2020multiplicity} showed that counterfactuals for one model can fail on other
models that are equally accurate. Upadhyay~\emph{et al.}~\cite{upadhyay2021robust} proposed computing recourse that
remains valid under a bounded set of model shifts. We take a simple, model-agnostic route suited to periodic
retraining: an option must also work on several bootstrap-retrained models, and its robustness is measured on
further retrained models held out from the search.

\subsection{Fairness in Credit Decisions}
Group fairness criteria compare decisions across groups, for example the rate of positive decisions (demographic
parity) or the rate among applicants who would repay (equal opportunity~\cite{hardt2016equality}). A common response
to a disparity is to remove the sensitive attribute from the model. Dwork~\emph{et al.}~\cite{dwork2012awareness}
argued that this ``fairness through unawareness'' is not sufficient, because other features can encode the same
information. We test this directly for age on our data, and keep the parts of the age policy that are about
eligibility rather than risk outside the model, as explicit rules.
